\subsection{Aufgabe 1 - Brennweite des Achromats}
Aus \cref{eq:f} folgt
\begin{equation}
    f=\frac{bg}{b+g}
\end{equation}
Die Fehlerformel bestimmt sich zu:
\begin{equation}
    \Delta f^{2}=\left(\left(\frac{-bg}{(b+g)^{2}}+\frac{g}{b+g}\right)\Delta b\right)^{2}+\left(\left(\frac{-bg}{(b+g)^{2}}+\frac{b}{b+g}\right)\Delta g\right)^{2}
\end{equation}

Berechnet wurde aus 7 messbaren Wertepaaren für b und g, jeweils 3 für I \& III und einem Wertepaar für II, sowie einem Ablesefehler \(\Delta b , \Delta g = \qty{0,1}{\cm}\):\\
\makebox[\textwidth][c]{\underline{\(f=\qty{9,93 \pm 0,05}{\cm}\)}}
Die richtige Brennweite von \(f=\qty{10}{\cm}\) liegt nicht innerhalb der wahrscheinlich zu niedrig eingeschätzten Fehlergrenzen. Für erhöhte Fehler \(\Delta g, \Delta b = \qty{0,2}{\cm}\) ergibt sich:\\
\makebox[\textwidth][c]{\underline{\(f=\qty{9,93 \pm 0,08}{\cm}\)}}

Grafisch ermittelt wurden in \cref{fig:Brennweite.pdf}:\\
\makebox[\textwidth][c]{\underline{\(f=\qty{10 \pm 0,6}{\cm}\)}}
In \cref{fig:Strahlengang.pdf} wurden mit gegebener Brennweite \(f=\qty{10}{\cm}\) und den gemessenen Gegenstandsweiten \(g\) grafisch \(b\)\&\(B\) ermittelt. Aufgrund zeichnerischen und ableserischen Ungenauigkeiten wurde ein Fehler von \(\Delta b,\Delta B=\qty{0,3}{\cm}\) geschätzt. 
Die ermittelten Ergebnisse lagen im angegebenen Fehlerbereich der im Versuch gemessenen Werte (siehe Tabelle in \cref{fig:Strahlengang.pdf}).\\\\
Als dritten Teil der Aufgabe 1 folgt die Auswertungstabelle:
\begin{table}[htbp]
\centering
\renewcommand{\arraystretch}{1.5}
\begin{tabularx}{0.8\textwidth}{c TTcTc}
\toprule
\textbf{Nr.} & \textbf{g} & \textbf{b} & \textbf{Art} & \textbf{Richtung} & \boldmath\(\beta\) \\
\midrule
I & \(\infty > g > 2f\) & \(f < b < 2f\) & reell & umgekehrt & \(<1\) \\
II & \(g = 2f\) & \(b  =2f\) & reell & umgekehrt & 1 \\
III & \(2f > g > f\) & \(b > 2f\) & reell & umgekehrt & \(>1\) \\
IV & \(g = f\) & \(b \rightarrow \infty\) & virtuell & N/A & N/A \\
V & \(f > g\) & N/A & virtuell & aufrecht & N/A \\
\bottomrule
\end{tabularx}
\caption{Messdaten Aufgabe 1}
\end{table}
\clearpage

\xfigpdf{p}{width=\textwidth}{Strahlengang.pdf}{}{Grafische Ermittlung der Brennweite \(f\)}
\clearpage
\xfigpdf{p}{width=\textwidth}{Brennweite.pdf}{}{Grafische Ermittlung von \(b\)}
\FloatBarrier

\subsection{Aufgabe 2 - Besselverfahren}
Mithilfe der \cref{eq:Bessel} wird die Brennweite einer neuen, bikonvexen Linse bei weißem Licht bestimmt. Dazu wurde bei einem Abstand \(l=\qty{65}{\cm} > 5f\) jeweils 3-mal ein Messwert für Schärfeposition \(x_1\) und \(x_2\) erhoben.
Wie erwartet ist der mittlere Fehler des Mittelwerts \(\sigma_{\bar{f}}\) aufgrund der sehr kleinen Anzahl an erhobenen Messungen sehr klein und deswegen im Vergleich zu den Ablesefehlern von \(\Delta x_{1},\Delta x_{2}=\qty{0,1}{\cm}\) zu vernachlässigen.

\begin{table}[h]
\centering
\begin{tabularx}{\textwidth}{cZZZZZZZZ}
\toprule
Nr. & l\;[\unit{\cm}] & x_1\;[\unit{\cm}] & x_2\;[\unit{\cm}] & d\;[\unit{\cm}] & \bar{d}\;[\unit{\cm}] & f\;[\unit{\cm}] & \bar{f}\;[\unit{\cm}] & \sigma_{\bar{f}}\;[\unit{\cm}]\\
\midrule 
 1 & 65     & 16.3    & 49.2 & 32.9 & 33.00   & 12.09 & 12.07 & 0.01 \\ 
 2 &        & 16.2   & 49.2 & 33     &         & 12.06  &      & \\ 
 3 &        & 16.4   & 49.4 & 33     &         & 12.06  &      & \\ 
\bottomrule
\end{tabularx}
% \label{}
\caption{Messdaten Besselverfahren}
\end{table}

Aus \(\Delta d=\sqrt{(-\Delta x_1)^{2}+(\Delta x_2)^{2}}=\sqrt{(-\qty{0,1}{\cm})^{2}+(\qty{0,1}{\cm})^{2}}=\qty{0,14}{\cm}\)
sowie \(\Delta l=\qty{0,2}{\cm}\) ergibt sich mit 
\begin{equation}
    \Delta f=\sqrt{\left(\left(\frac{1}{2}-\frac{-d^2+l^2}{4l^2}\right) \Delta l\right)^2+\left(\frac{-d}{2l} \Delta d\right)^2}
\end{equation}
\makebox[\textwidth][c]{\underline{\(f_{\text{weiß}}=\qty{12,07 \pm 0,07}{\cm}\)}}

\subsection{Aufgabe 3 - Besselverfahren mit rot und blau}
Wie in Aufgabe 2 werden dieselben Rechnungen gemacht. Es ergibt sich bei rot ein niedrigerer Wert für \(d=\qty{32,7}{\cm}\), bei blau ein höherer Wert von \(d=\qty{33,5}{\cm}\).\\
Damit ist:\\
\makebox[\textwidth][c]{\underline{\(f_{\text{rot}}=\qty{12,15 \pm 0,07}{\cm}\)}}\\
\makebox[\textwidth][c]{\underline{\(f_{\text{blau}}=\qty{11,95 \pm 0,07}{\cm}\)}}\\
Die Fehler unterscheiden sich nur marginal an nicht signifikanten Stellen.

\subsection{Aufgabe 4 - Sphärische Abberation}
\paragraph{Lochblende:}
Bei der Lochblende war das Bild wie zu erwarten schärfer, das Licht konnte nur im Zentrum
auf die Linse treffen kann, dadurch verringern sich die Effekte der sphärischen Abberation. Der Abstand auf \(d=\qty{32}{\cm}\) hat sich verringert, damit steigt die Brennweite nach der Besselformel \cref{eq:Bessel}. 
\paragraph{Ringblende:}
Das Bild war nicht mehr so scharf zu erkennen wie bei der Lochblende. Der Abstand hat sich auf \(d=\qty{33,4}{\cm}\) vergrößert, damit steigt \(f\) nach \cref{eq:Bessel}.

\subsection{Aufgabe 5 - Gitterkonstante}
Auf der optischen Bank wurde ein Gitter so positioniert, dass der Abstand \(b=\qty{25}{\cm}\), \(\Delta b=\qty{0,2}{\cm}\) des Zwischenbildes zur Objektivlinse mit \(f_1=\qty{4}{\cm}\) (nicht fehlerbehaftet) betrug. Das Zwischenbild wurde danach näher durch die Okularlinse untersucht.
Über folgende Gleichung kann der Abbildungsmaßstab \(\beta\), \(\Delta \beta\) im Zwischenbild bestimmt werden:
\begin{align}
    \beta = \frac{b}{f} - 1 && \Delta \beta = \frac{\Delta b}{f} && \beta=\num{5.50(5)}
\end{align}
Somit beträgt nach Umstellen der \cref{eq:beta} und dem Ablesen von \(B=\qty{0,05 \pm 0,005}{\cm}\) am Dia mit mm-Einteilung der Gitterabstand \(G\) als Gegenstandsgröße, sowie dem Fehler: 
\begin{equation}
    \Delta G = \sqrt{\left( \frac{1}{\beta} \Delta B \right)^2+ \left( -\frac{B}{\beta^2} \Delta \beta \right)^2}=\qty{9,12 e-6}{\m}
\end{equation}
\makebox[\textwidth][c]{\underline{\(G=\frac{B}{\beta}=\qty{90 \pm 9}{\micro\m}\)}}

\subsection{Aufgabe 6 - Auflösung des Mikroskops}
Ein verstellbarer Spalt mit Spaltbreite \(d=\qty{0,2}{\mm}\), \(\Delta d=\qty{0,1}{\mm}\) wurde im Abstand \(x=\qty{37}{\mm}\), \(\Delta x=\qty{2}{\mm}\) vom Gitter entfernt positioniert. \(d\) wurde solange verringert, bis durch Beugungseffekte die Gitterstrukturen durch die Okularlinse nicht mehr sichtbar waren.\\
Damit ist der Öffnungswinkel des Objektives (die Objektivlinse stand direkt hinter dem Spalt) unter Ausnutzung der Kleinwinkelnäherung folgend zu berechnen:
\begin{align}
    \alpha=\text{arctan}\left(\frac{\nicefrac{d}{2}}{x}\right)\approx\frac{\nicefrac{d}{2}}{x}
    \label{eq:alpha}
    && \Delta \alpha=\sqrt{\left(\frac{1}{2x}\Delta d\right)^{2}+\left(-\frac{\nicefrac{d}{2}}{x^{2}}\Delta x\right)^{2}}
\end{align}
\makebox[\textwidth][c]{\underline{\(\alpha=\qty{0,0027 \pm 0,0014}{\degree}\)}}

Nach \cref{eq:NA} kann mit den obigen Größen und der Wellenlänge \(\lambda=\qty{550}{\nm}\) sowie Brechungsindex \(n=1\) für Luft der minimale Abstand \(G_{\text{min}}\) zweier Gegenstandspunkte dieses Mikroskopaufbaus berechnet werden.
Bei einem solchen kleinen Winkel ist auch für \(G_{\text{min}}\) die Kleinwinkelnäherung nutzbar. 
\begin{align}
    G_{\text{min}}=\frac{0,61\lambda}{n\cdot\alpha}&&
    \label{GMIN}
    \Delta G_{\text{min}}=\sqrt{\left(\frac{-0,61\lambda}{n\cdot\alpha^{2}}\Delta \alpha\right)^{2}}
\end{align}
Damit beträgt der durch das Beobachten des Zwischenbildes ermittelte Gitterabstand:\\
\makebox[\textwidth][c]{\underline{\(G_{\text{min}}=\qty{124 \pm 64}{\micro\m}\)}}

Bei blauem Licht hat sich ein Spaltabstand von \(d=\qty{0,3 \pm 0,1}{\mm}\) ergeben und bei rotem Licht von \(d=\qty{0,2 \pm 0,1}{\mm}\). Also tendenziell bedeutet eine höhere Wellenlänge eine größere Spaltbreite, dies wiederum führt durch \(G_{\text{min}}\sim\frac{1}{d}\) (siehe \cref{eq:alpha,GMIN}) zu einem kleineren minimal möglichen Gegenstandsabstand. Dies macht auch Sinn, da Beugung erst bei Strukturen in der Größenordnungen der bestrahlenden Wellenlänge auftritt.